\documentclass[12pt]{ximera}
\typeout{Start loading xmPreamble.tex}%

% Add here extra macro's that are loaded automatically by all documents of claas 'ximera' or 'xourse' in this repo

%%
%%  Example:
%%
% \newcommand{\R}{\mathbb{R}}
\usepackage{tikz}
\usetikzlibrary{positioning, arrows.meta}
\usepackage{multicol}
\usepackage{enumitem}
\usepackage{caption}
\usepackage{amsmath}
\usepackage{amsthm}
\usepackage{fontenc}

%% ---------------------------------------------------------------
%% Shared worksheet macros.
%%
%% These came from the Summer 2026 worksheet set, where each worksheet carried its
%% own copy. They have to live here instead: a xourse inlines every activity into a
%% single document, so duplicate \newcommand definitions across activities collide,
%% and the xourse gobbles \input so a per-topic macro file never loads.
%%
%% They use \providecommand, NOT \newcommand. ximera.cls loads this file automatically,
%% and every activity also carries an explicit \typeout{Start loading xmPreamble.tex}%

% Add here extra macro's that are loaded automatically by all documents of claas 'ximera' or 'xourse' in this repo

%%
%%  Example:
%%
% \newcommand{\R}{\mathbb{R}}
\usepackage{tikz}
\usetikzlibrary{positioning, arrows.meta}
\usepackage{multicol}
\usepackage{enumitem}
\usepackage{caption}
\usepackage{amsmath}
\usepackage{amsthm}
\usepackage{fontenc}

%% ---------------------------------------------------------------
%% Shared worksheet macros.
%%
%% These came from the Summer 2026 worksheet set, where each worksheet carried its
%% own copy. They have to live here instead: a xourse inlines every activity into a
%% single document, so duplicate \newcommand definitions across activities collide,
%% and the xourse gobbles \input so a per-topic macro file never loads.
%%
%% They use \providecommand, NOT \newcommand. ximera.cls loads this file automatically,
%% and every activity also carries an explicit \typeout{Start loading xmPreamble.tex}%

% Add here extra macro's that are loaded automatically by all documents of claas 'ximera' or 'xourse' in this repo

%%
%%  Example:
%%
% \newcommand{\R}{\mathbb{R}}
\usepackage{tikz}
\usetikzlibrary{positioning, arrows.meta}
\usepackage{multicol}
\usepackage{enumitem}
\usepackage{caption}
\usepackage{amsmath}
\usepackage{amsthm}
\usepackage{fontenc}

%% ---------------------------------------------------------------
%% Shared worksheet macros.
%%
%% These came from the Summer 2026 worksheet set, where each worksheet carried its
%% own copy. They have to live here instead: a xourse inlines every activity into a
%% single document, so duplicate \newcommand definitions across activities collide,
%% and the xourse gobbles \input so a per-topic macro file never loads.
%%
%% They use \providecommand, NOT \newcommand. ximera.cls loads this file automatically,
%% and every activity also carries an explicit \input{../xmPreamble}, so this file is
%% read twice. That was harmless while it held only \usepackage lines (LaTeX skips a
%% package it has already loaded) but a second \newcommand is a hard error.
%%
%%   \blank[width]        fill-in-the-blank rule (default 2.5cm)
%%   \showwork            "show and explain your work" reminder
%%   \showworkline        ... plus which schedule line was used
%%   \showworkyear        ... plus which year's schedule was used
%%   \schedhead{...}      centred bold caption above a tiered-rate table
%%   \samesquares[scale]{left}{right}   two equal-area squares, labelled
%%   \samecubes[scale]{left}{right}     two equal-volume cubes, labelled
%% ---------------------------------------------------------------

\providecommand{\blank}[1][2.5cm]{\underline{\hspace{#1}}}
\providecommand{\showwork}{\textbf{REMEMBER TO SHOW AND EXPLAIN YOUR WORK.}}
\providecommand{\showworkline}{\textbf{REMEMBER TO SHOW AND EXPLAIN YOUR WORK.
  Explanation should include how and why you know which line of the
  schedule was used to calculate this amount.}}
\providecommand{\showworkyear}{\begin{sloppypar}\textbf{REMEMBER TO SHOW AND
  EXPLAIN YOUR WORK. Explanation should include how and why you know which
  line of the year's tax schedule was used to calculate this tax
  amount.}\end{sloppypar}}
\providecommand{\schedhead}[1]{\begin{center}\textbf{#1}\end{center}}

\providecommand{\samesquares}[3][1]{%
\begin{center}
{\footnotesize These squares are the \textbf{same size}.}

\vspace{1.5mm}
\begin{tikzpicture}[scale=#1,every node/.style={scale=#1}]
  \draw (0,0) rectangle (2.4,2.4);
  \node[anchor=north west] at (0.15,2.25) {Area:};
  \node[anchor=east]  at (-0.15,1.2) {#2};
  \node[anchor=north] at (1.2,-0.15) {#2};
  \begin{scope}[xshift=4.6cm]
    \draw (0,0) rectangle (2.4,2.4);
    \node[anchor=north west] at (0.15,2.25) {Area:};
    \node[anchor=east]  at (-0.15,1.2) {#3};
    \node[anchor=north] at (1.2,-0.15) {#3};
  \end{scope}
\end{tikzpicture}
\end{center}}

\providecommand{\samecubes}[3][1]{%
\begin{center}
{\footnotesize These cubes are the \textbf{same size}.}

\vspace{1.5mm}
\begin{tikzpicture}[scale=#1,every node/.style={scale=#1}]
  \foreach \dx in {0,5.2} {
    \begin{scope}[xshift=\dx cm]
      \draw (0,0) rectangle (2.0,2.0);
      \draw (0,2.0) -- (0.65,2.6) -- (2.65,2.6) -- (2.0,2.0);
      \fill[black!12] (2.0,0) -- (2.65,0.6) -- (2.65,2.6) -- (2.0,2.0) -- cycle;
      \draw (2.0,0) -- (2.65,0.6) -- (2.65,2.6) -- (2.0,2.0);
      \node[anchor=north west] at (0.1,1.9) {Volume:};
    \end{scope}
  }
  \node[anchor=east]  at (-0.15,1.0) {#2};
  \node[anchor=north] at (1.0,-0.15) {#2};
  \node[anchor=west]  at (2.25,0.4) {#2};
  \node[anchor=east]  at (5.05,1.0) {#3};
  \node[anchor=north] at (6.2,-0.15) {#3};
  \node[anchor=west]  at (7.45,0.4) {#3};
\end{tikzpicture}
\end{center}}

%% NOTE: do NOT add \usepackage to individual activity .tex files.
%% When a xourse inlines an activity it gobbles \documentclass and \input, but a bare
%% \usepackage still executes -- after \begin{document} -- and errors out.
%% All shared packages and macros belong here.
%%
%% ximera.cls already loads: amsmath, amssymb, amsthm, cancel, enumitem, graphicx,
%% hyperref, listings, pgfplots, tikz, url, xcolor. No need to re-load those.
, so this file is
%% read twice. That was harmless while it held only \usepackage lines (LaTeX skips a
%% package it has already loaded) but a second \newcommand is a hard error.
%%
%%   \blank[width]        fill-in-the-blank rule (default 2.5cm)
%%   \showwork            "show and explain your work" reminder
%%   \showworkline        ... plus which schedule line was used
%%   \showworkyear        ... plus which year's schedule was used
%%   \schedhead{...}      centred bold caption above a tiered-rate table
%%   \samesquares[scale]{left}{right}   two equal-area squares, labelled
%%   \samecubes[scale]{left}{right}     two equal-volume cubes, labelled
%% ---------------------------------------------------------------

\providecommand{\blank}[1][2.5cm]{\underline{\hspace{#1}}}
\providecommand{\showwork}{\textbf{REMEMBER TO SHOW AND EXPLAIN YOUR WORK.}}
\providecommand{\showworkline}{\textbf{REMEMBER TO SHOW AND EXPLAIN YOUR WORK.
  Explanation should include how and why you know which line of the
  schedule was used to calculate this amount.}}
\providecommand{\showworkyear}{\begin{sloppypar}\textbf{REMEMBER TO SHOW AND
  EXPLAIN YOUR WORK. Explanation should include how and why you know which
  line of the year's tax schedule was used to calculate this tax
  amount.}\end{sloppypar}}
\providecommand{\schedhead}[1]{\begin{center}\textbf{#1}\end{center}}

\providecommand{\samesquares}[3][1]{%
\begin{center}
{\footnotesize These squares are the \textbf{same size}.}

\vspace{1.5mm}
\begin{tikzpicture}[scale=#1,every node/.style={scale=#1}]
  \draw (0,0) rectangle (2.4,2.4);
  \node[anchor=north west] at (0.15,2.25) {Area:};
  \node[anchor=east]  at (-0.15,1.2) {#2};
  \node[anchor=north] at (1.2,-0.15) {#2};
  \begin{scope}[xshift=4.6cm]
    \draw (0,0) rectangle (2.4,2.4);
    \node[anchor=north west] at (0.15,2.25) {Area:};
    \node[anchor=east]  at (-0.15,1.2) {#3};
    \node[anchor=north] at (1.2,-0.15) {#3};
  \end{scope}
\end{tikzpicture}
\end{center}}

\providecommand{\samecubes}[3][1]{%
\begin{center}
{\footnotesize These cubes are the \textbf{same size}.}

\vspace{1.5mm}
\begin{tikzpicture}[scale=#1,every node/.style={scale=#1}]
  \foreach \dx in {0,5.2} {
    \begin{scope}[xshift=\dx cm]
      \draw (0,0) rectangle (2.0,2.0);
      \draw (0,2.0) -- (0.65,2.6) -- (2.65,2.6) -- (2.0,2.0);
      \fill[black!12] (2.0,0) -- (2.65,0.6) -- (2.65,2.6) -- (2.0,2.0) -- cycle;
      \draw (2.0,0) -- (2.65,0.6) -- (2.65,2.6) -- (2.0,2.0);
      \node[anchor=north west] at (0.1,1.9) {Volume:};
    \end{scope}
  }
  \node[anchor=east]  at (-0.15,1.0) {#2};
  \node[anchor=north] at (1.0,-0.15) {#2};
  \node[anchor=west]  at (2.25,0.4) {#2};
  \node[anchor=east]  at (5.05,1.0) {#3};
  \node[anchor=north] at (6.2,-0.15) {#3};
  \node[anchor=west]  at (7.45,0.4) {#3};
\end{tikzpicture}
\end{center}}

%% NOTE: do NOT add \usepackage to individual activity .tex files.
%% When a xourse inlines an activity it gobbles \documentclass and \input, but a bare
%% \usepackage still executes -- after \begin{document} -- and errors out.
%% All shared packages and macros belong here.
%%
%% ximera.cls already loads: amsmath, amssymb, amsthm, cancel, enumitem, graphicx,
%% hyperref, listings, pgfplots, tikz, url, xcolor. No need to re-load those.
, so this file is
%% read twice. That was harmless while it held only \usepackage lines (LaTeX skips a
%% package it has already loaded) but a second \newcommand is a hard error.
%%
%%   \blank[width]        fill-in-the-blank rule (default 2.5cm)
%%   \showwork            "show and explain your work" reminder
%%   \showworkline        ... plus which schedule line was used
%%   \showworkyear        ... plus which year's schedule was used
%%   \schedhead{...}      centred bold caption above a tiered-rate table
%%   \samesquares[scale]{left}{right}   two equal-area squares, labelled
%%   \samecubes[scale]{left}{right}     two equal-volume cubes, labelled
%% ---------------------------------------------------------------

\providecommand{\blank}[1][2.5cm]{\underline{\hspace{#1}}}
\providecommand{\showwork}{\textbf{REMEMBER TO SHOW AND EXPLAIN YOUR WORK.}}
\providecommand{\showworkline}{\textbf{REMEMBER TO SHOW AND EXPLAIN YOUR WORK.
  Explanation should include how and why you know which line of the
  schedule was used to calculate this amount.}}
\providecommand{\showworkyear}{\begin{sloppypar}\textbf{REMEMBER TO SHOW AND
  EXPLAIN YOUR WORK. Explanation should include how and why you know which
  line of the year's tax schedule was used to calculate this tax
  amount.}\end{sloppypar}}
\providecommand{\schedhead}[1]{\begin{center}\textbf{#1}\end{center}}

\providecommand{\samesquares}[3][1]{%
\begin{center}
{\footnotesize These squares are the \textbf{same size}.}

\vspace{1.5mm}
\begin{tikzpicture}[scale=#1,every node/.style={scale=#1}]
  \draw (0,0) rectangle (2.4,2.4);
  \node[anchor=north west] at (0.15,2.25) {Area:};
  \node[anchor=east]  at (-0.15,1.2) {#2};
  \node[anchor=north] at (1.2,-0.15) {#2};
  \begin{scope}[xshift=4.6cm]
    \draw (0,0) rectangle (2.4,2.4);
    \node[anchor=north west] at (0.15,2.25) {Area:};
    \node[anchor=east]  at (-0.15,1.2) {#3};
    \node[anchor=north] at (1.2,-0.15) {#3};
  \end{scope}
\end{tikzpicture}
\end{center}}

\providecommand{\samecubes}[3][1]{%
\begin{center}
{\footnotesize These cubes are the \textbf{same size}.}

\vspace{1.5mm}
\begin{tikzpicture}[scale=#1,every node/.style={scale=#1}]
  \foreach \dx in {0,5.2} {
    \begin{scope}[xshift=\dx cm]
      \draw (0,0) rectangle (2.0,2.0);
      \draw (0,2.0) -- (0.65,2.6) -- (2.65,2.6) -- (2.0,2.0);
      \fill[black!12] (2.0,0) -- (2.65,0.6) -- (2.65,2.6) -- (2.0,2.0) -- cycle;
      \draw (2.0,0) -- (2.65,0.6) -- (2.65,2.6) -- (2.0,2.0);
      \node[anchor=north west] at (0.1,1.9) {Volume:};
    \end{scope}
  }
  \node[anchor=east]  at (-0.15,1.0) {#2};
  \node[anchor=north] at (1.0,-0.15) {#2};
  \node[anchor=west]  at (2.25,0.4) {#2};
  \node[anchor=east]  at (5.05,1.0) {#3};
  \node[anchor=north] at (6.2,-0.15) {#3};
  \node[anchor=west]  at (7.45,0.4) {#3};
\end{tikzpicture}
\end{center}}

%% NOTE: do NOT add \usepackage to individual activity .tex files.
%% When a xourse inlines an activity it gobbles \documentclass and \input, but a bare
%% \usepackage still executes -- after \begin{document} -- and errors out.
%% All shared packages and macros belong here.
%%
%% ximera.cls already loads: amsmath, amssymb, amsthm, cancel, enumitem, graphicx,
%% hyperref, listings, pgfplots, tikz, url, xcolor. No need to re-load those.

\title{Introduction: Connecting the Math and the Context}
\begin{document}
\begin{abstract}
How to write explanations in this course: every answer should name something
\emph{mathematical} --- the $x$ and $y$ values you are reading off a graph --- and
something about the \emph{context} that graph represents. Practice on a
distance-from-home graph and on a jet stream temperature map.
\end{abstract}
\maketitle

\section*{Connect the ``Math'' and the ``Context''}

In this class, the \textbf{explanation} behind your answers is just as important or
more important than the answer itself. That explanation should be connecting the
\textbf{mathematics} of the problem or information to the \textbf{context} which the
problem or information is related to in a ``real world'' setting.

\subsection*{Problem Set 1: Shana's Distance from Home}

In your answers to the following questions, write in \textbf{complete sentences}.
State \textbf{explicitly} what you see in the \textbf{graph} and why that led you to
your answer choice. This means each of your answers should mention something
\textbf{mathematical}, i.e.\ what the $x$ and $y$ values from the graph are,
\textbf{AND} something about the \textbf{context}, i.e.\ something about what this
graph is representing in the ``real world.''

\begin{center}
\includegraphics[width=0.85\textwidth]{introduction/shana_graph.png}
\end{center}

Above is a graph which gives information about Shana's distance from home throughout
the day. Use this graph to answer the following questions.

\begin{problem}
What does the vertical axis represent? Choose a specific point, like $(6,0)$ to help
describe this (don't use $(6,0)$)!

\begin{freeResponse}
\end{freeResponse}

\begin{solution}
The vertical axis represents the distance Shana is from her home, given in miles. For
example, the point $(11, 10.25)$ is on the graph, meaning at $x = 11$, $y = 10.25$,
meaning Shana is 10.25 miles from her home.

\textbf{Note:} It is important to mention what $y$ equals, and then state what that
means in the context (i.e.\ something about Shana and her distance from home).
\end{solution}
\end{problem}

\begin{problem}
What does the horizontal axis represent? Choose a specific point, like $(6,0)$ to help
describe this (don't use $(6,0)$)!

\begin{freeResponse}
\end{freeResponse}

\begin{solution}
The horizontal axis represents the time in hours since the start of the day
(midnight). For example, the point $(11, 10.25)$ is on the graph, meaning at $x = 11$,
$y = 10.25$, meaning at 11~AM, Shana is 10.25 miles from her home.

\textbf{Note:} It is important to mention what $x$ equals, and then state what that
means in the context (i.e.\ something about what time it is and/or how many hours have
passed).
\end{solution}
\end{problem}

\begin{problem}
Why are there no negative vertical values? Use a specific value as an example in your
answer.

\begin{freeResponse}
\end{freeResponse}

\begin{solution}
There are no negative values because the vertical axis is measuring distance, and we
do not describe distance as negative. So ``$-4$ miles'' does not have meaning, it
would just be ``4'' in the graph.
\end{solution}
\end{problem}

\begin{problem}
What can you say about when Shana woke up? Why isn't this asking you ``When did Shana
wake up?''

\begin{freeResponse}
\end{freeResponse}

\begin{solution}
The $y$-value is not greater than 0 until $x = 7$. This means Shana's distance from
home does not change from 0 miles until 7 hours (or 7~AM). This means that we can say
Shana got up before 7~AM, though we cannot be more specific than that.

Since we cannot know exactly what time she woke up, that is why the question does not
ask for that specificity.

\textbf{Note:} We must talk about the ``$x$'' and ``$y$'' values from the graph in
addition to the 0 miles and 7 hours/7~AM, as this is connecting the math ($x$ \& $y$
from graph) to the context (Shana's distance at a certain time).
\end{solution}
\end{problem}

\begin{problem}
Did Shana likely have lunch at work?

\begin{freeResponse}
\end{freeResponse}

\begin{solution}
Shana likely did NOT have lunch at work. Between the $x$ values of 12 and 13, the
$y$-value goes down a little bit. The $x$ value of 12 and 13 is about 12~PM and 1~PM,
which is a reasonable lunch time. Since the distance from home changes during this
time, we can assume that Shana moved away from work, and since this is around lunch
time, we can assume she left work for lunch.

\textbf{Note:} We must talk about the ``$x$'' and ``$y$'' values from the graph in
addition to the miles/distance changing and the 12 and 13 hours (12 \& 1~PM). This is
connecting the math ($x$ \& $y$ from graph) to the context (Shana's distance at a
certain time).
\end{solution}
\end{problem}

\begin{problem}
What was Shana's average speed getting to work, in miles per hour?

Shana's average speed getting to work was $\answer[tolerance=0.5]{10.25}$ mph.

\begin{freeResponse}
\end{freeResponse}

\begin{solution}
Speed is distance divided by time. So we need to use the graph to find out how far
Shana went from home to work and then divide it by the amount of time it took her to
go that distance.

To find Shana's distance from home to work, we look at $(7,0)$ which is when the
$y$-value stops being 0 and $(8, 10.25)$ which is when the $y$-value starts to stay
the same again. Since the $y$-values are the distance, we can do $10.25 - 0 = 10.25$
miles is the difference, so Shana traveled 10.25 miles.

To find the time from home to work, we look at the same points: $(7,0)$ and
$(8,10.25)$ but now at the $x$-values as the $x$-values tell us the time. We can do
$8 - 7 = 1$ to see that it took Shana 1 hour to travel that distance.

A distance of 10.25 miles divided by 1 hour means she traveled at 10.25 miles per hour
from home to work.

\textbf{Note:} As before, you must reference the $x$ \& $y$ values (and/or points) from
the graph and how they relate to time and distance in order to connect the math (the
graph) to the context (Shana's distance at a time).
\end{solution}
\end{problem}

\begin{problem}
Would the graph being steeper mean Shana's speed would be faster or slower? Explain how
you know.

\begin{freeResponse}
\end{freeResponse}

\begin{solution}
Since speed is distance divided by time, something steeper would mean moving faster.
This is because something steep would mean the $y$-values, distance, would have a
bigger (or same)* difference, which means a longer (or same) distance, while the
$x$-values, time, would have a smaller (or same) difference, meaning less time (or same
time) for the travel. This would mean we would be dividing a bigger number by a smaller
number than before, making the result, which would mean the speed, bigger.

* Only one of the values actually needs to change for the line to be steeper: either
the distance increases (bigger difference in $y$-values) or the time decreases (smaller
difference in $x$-values), or both can happen to make things steeper! As long as you
mention at least one of them changing specifically, and how that is affecting the
result (speed), that is acceptable for this question.

\textbf{Note:} You must talk about the change in the $x$ \& $y$ values changing and how
that connects to the distance and time in order to connect the math and the context.
You should also reference the operation done with these values (dividing) as this is
important to understanding how the result (speed) is affected by the $x$ \& $y$
changes.
\end{solution}
\end{problem}

\begin{problem}
Explain why we cannot have two points which have the same horizontal value in this
graph. Make up two specific points which would have the same horizontal values in the
graph to use as an example in your answer.

\begin{freeResponse}
\end{freeResponse}

\begin{solution}
Two points having the same horizontal value would mean something like the points
$(11, 10.25)$ and $(11, 7.5)$ being part of the graph. This would mean since $x = 11$
for both, that after 11 hours (at 11~AM) Shana was BOTH 10.25 miles away AND 7.5 miles
away, since those are the $y$-values. Since it is not possible to be in two places at
once, this is not possible, which is why two points cannot have the same horizontal
value.
\end{solution}
\end{problem}

\begin{problem}
What does it mean if two points have the same vertical value on the graph? Choose two
specific points on the graph with the same vertical values to use as an example in your
answer.

\begin{freeResponse}
\end{freeResponse}

\begin{solution}
Two points which have the same vertical values could be the points $(11, 10.25)$ and
$(14, 10.25)$ from the graph. This means at both $x = 11$ and $x = 14$, meaning at 11
hours (11~AM) and 14 hours (2~PM), the $y$-value was 10.25, meaning Shana was 10.25
miles away from home at both of those times. Generally, having the same $y$-value for
two points means Shana was the same distance from home (in miles) at those two
different times (different $x$-values).
\end{solution}
\end{problem}

\begin{problem}
It is reasonable to assume Shana was at work and not moving during the 8~AM to 12~PM
time frame. However, this does not have to be the case. Describe a scenario for when
Shana's distance from home is fixed, but Shana is actually moving.

\begin{freeResponse}
\end{freeResponse}

\begin{solution}
If Shana is a fixed distance from home, but moving, she could be moving in a circle
around her house:

\begin{center}
\includegraphics[width=0.4\textwidth]{introduction/circle_house.png}
\end{center}

By moving in a circle around her house, Shana would be a fixed distance from her home,
as the radius of the circle would not be changing, but she would be moving along the
edge of the circle and thus moving in relation to her house.
\end{solution}
\end{problem}

\begin{problem}
After work, Shana went out to dinner and then to the grocery store after. Can we tell
if Shana went toward or away from her work to get dinner and go to the grocery store?
Explain why or why not.

\begin{freeResponse}
\end{freeResponse}

\begin{solution}
No, we cannot tell if Shana went toward or away from work when she left for dinner and
the grocery store. This is because the distance given by the $y$-values are distance
from her home, not in relation to anything else. So when Shana's distance increases
from home, we don't know what direction she is going from her home. The direction could
be the same direction as another time she left home, like north, or it could be in a
different direction, like south or west.

\begin{center}
\includegraphics[width=0.5\textwidth]{introduction/house_distances.png}
\end{center}

Each of these distances would have $y$-values of 10, 5, or 3 in the graph, but are in
different directions from Shana's home.
\end{solution}
\end{problem}

\subsection*{Problem Set 2: Jet Streams and Temperature}

The graph below is meant to represent what locations in the U.S.\ share a common
general temperature in a given season (temperatures are in $^\circ$F) in order to
illustrate the idea of jet streams and their general path by how they affect the
temperature in the U.S. However, that information is not particularly important for our
purposes. Use this graph to answer the following questions.

\begin{center}
\includegraphics[width=0.85\textwidth]{introduction/jetstream_map.png}

\url{https://www.weather.gov/jetstream/ll_analyze_temp}
\end{center}

\begin{problem}
Locate each of the following cities on the map, and use its position relative to the
jet stream lines to estimate a temperature for it. You do not need to provide a written
explanation for your temperature choices for this problem.

\begin{align*}
\text{Columbus, OH } &\answer[tolerance=10]{48}\,^\circ\text{F}
&\text{Seattle, WA } &\answer[tolerance=10]{40}\,^\circ\text{F} \\
\text{Los Angeles, CA } &\answer[tolerance=10]{42}\,^\circ\text{F}
&\text{Dallas, TX } &\answer[tolerance=10]{55}\,^\circ\text{F} \\
\text{New Orleans, LA } &\answer[tolerance=10]{80}\,^\circ\text{F}
&\text{Rapid City, SD } &\answer[tolerance=10]{10}\,^\circ\text{F} \\
\text{Nashville, TN } &\answer[tolerance=10]{70}\,^\circ\text{F}
&\text{Miami, FL } &\answer[tolerance=10]{95}\,^\circ\text{F}
\end{align*}

\begin{solution}
Reading each city's approximate position against the labelled jet stream lines gives
roughly: Columbus, OH $48^\circ$; Seattle, WA $40^\circ$; Los Angeles, CA $42^\circ$;
Dallas, TX $55^\circ$; New Orleans, LA $80^\circ$; Rapid City, SD $10^\circ$;
Nashville, TN $70^\circ$; Miami, FL $95^\circ$.

These are estimates read off a map, so a range of nearby answers is reasonable.
\end{solution}
\end{problem}

\begin{problem}
For which city above did you find it most difficult to decide on an estimated
temperature? Why?

\begin{freeResponse}
\end{freeResponse}

\begin{solution}
I found it most difficult to decide on a temperature for Rapid City, SD. I found it
difficult because it is almost in the dead center of the ``20 degree loop.'' So I
wasn't sure if that should be a temperature like 0 or 10 or if it would actually be
closer to 20 degrees. There was not another value to know it should fall between like
with other cities.

\textbf{Note:} Another city (and explanation with it) could also be acceptable here!
For example, I could see a good argument made for difficulty in choosing a temperature
choice for Dallas, TX or Miami, FL.
\end{solution}
\end{problem}

\begin{problem}
For which cities (one or two) above do you think there may be the widest variety in the
estimated temperature? Explain why you think this.

\begin{freeResponse}
\end{freeResponse}

\begin{solution}
I think there will be the most variation around Dallas, TX and Nashville, TN. This is
because both of those cities are in/around a cluster of degree lines, so those cities'
true locations would be very important for determining a temperature. Since these
cities are just in approximate areas, there could be a lot of variation in temperature
choices around those jet stream lines.
\end{solution}
\end{problem}

\begin{problem}
For which cities (one or two) above do you think there may be the least variety in the
estimated temperature? Explain why you think this.

\begin{freeResponse}
\end{freeResponse}

\begin{solution}
As a converse to the above, I think there may be a lot of agreement on Seattle, WA and
New Orleans, LA because they are so close/almost on the temperature lines generally,
that I think the temperature line's value will be consistently chosen for those cities.

\textbf{Note:} In the example answers above, the reasoning is given based on the map
information, and not based on personal thoughts about those cities, or thoughts on the
accuracy of those cities' marked locations on the map. The arguments are all made
assuming the map and its contents are correct and accurate. This is an important aspect
to remember when writing your own explanations and work for this course.
\end{solution}
\end{problem}

\end{document}
